\section[Local comparisons on horizontal tests of the split determinant]{Local comparisons on horizontal tests\\of the split determinant}\label{sec:split-local}

We retain the split imaginary field, the moving split primes, and the determinants and measures of Section~\ref{sec:split-measures}. In particular, for the reducible split-product input, \(E\) is non-CM, \(E(\mathbb Q)[2]\ne0\), \(K=\mathbb Q(\sqrt{-d})\) with \(d\equiv7\pmod {16}\) prime, and every prime of \(2N\) splits in \(K\). The primes \(r_i\), their limiting Tate matrix \(g\), and \(a_*=\operatorname{tr}g\ne2\) satisfy the conditions used in \eqref{eq:A2}--\eqref{eq:A3}. Until the rank argument at the center, no hypothesis about analytic rank is imposed.

We first compare local conditions after specializing to a high finite dyadic character. The resulting horizontal divisibility will be returned to the integral two-variable ring; the residual tests then make the quotient a unit. The center arguments are separated into a rank deduction from \(s_2(E/K)=1\) and an exact valuation calculation once the product has a simple zero. Finally, we isolate a characteristic-zero Heegner rank test whose proof uses only the derivative identities and reciprocity.

The local logarithm and inert derivative constructions also have a
variant for a fixed two-dimensional Hecke component of the Tate module
of a modular abelian variety.  We state the extra hypotheses where that
variant is used, with fixed denominators for its coefficient projector
and polarization.  The integral unit comparison in this section uses
the reducible residual calculations of Section~\ref{sec:split-measures}.

\subsection{Odd local conditions and projected dyadic Kummer lines}

Specialize \(C_r(t,u)\) at sufficiently high finite \(\theta\) in \(t\); put \(\Theta=\mathbb Z_2[\operatorname{im}\theta]\), \(A_0=\Theta[[u]]\), \(A=A_0[1/2]\), a regular one-dimensional UFD. Denote the limiting twisted coefficient by \(M\), with scalar \(\chi=\theta\Psi_u\) in stage/limit notation. Work on arithmetic diagrams over \(A\); on the measure side denote the analogously 2-inverted integer power-series ring using extended constants (with \(\mathcal V^*\) and fixed finite additional scalars as needed) by \(A'\). The conductor depth \(m\) of \eqref{eq:M2} is fixed throughout any one such varying-prime test. We require \(\theta\) of order greater than four on dyadic inertia and still nontrivial on inertia over any indicated fixed local extension.

\begin{lemma}[Odd local conditions]\label{split-local:odd}
At the places over the moving primes \(r\), the full local complexes are acyclic over \(A\). At each fixed odd place \(v\), the characteristic-zero unramified condition
\[
 U_v=[M^{I_v}\xrightarrow{\mathrm{Fr}_v-1}M^{I_v}]
 \qquad\text{in degrees }0,1
\]
contains local \(H^0\), injects into local \(H^1\), and is exactly complementary to the corresponding condition on the conjugate Tate dual. The product of the singular determinants at the two places over \(q\in S^\circ\) is, up to a unit,
\[
 D_q=P_q(\chi(\mathrm{Fr}_v))P_q(\chi(\mathrm{Fr}_v)^{-1}).
\]
These factors are nonzero in the present horizontal test.
\end{lemma}

\begin{proof}
At the places of \(r\), use the inertia/residue calculation on limits: a pro-2 tame generator acts with difference \(u\) (up to unit); where this vanishes on a residue field, the residue Frobenius and its singular twist have action-minus-one invertible by \(r_i\to1,\ \operatorname{tr}g\ne2\) and \(\theta(\mathrm{Fr}_R)\to1\). 

At each fixed odd place \(v\), replace the full condition by \(U_v\). Its quotient in the full local complex is represented, for cohomology and determinant purposes, by
\[
                     [\,M_{I_v}(-1)\longrightarrow M_{I_v}(-1)\,]
                         \qquad (1,2)
\]
with Frobenius differential. Indeed the scalar is unramified and the Tate inertia differential (after discarding the pro-(prime-to-2) kernel) has constant free invariant and coinvariant spaces after inverting 2. Relative to its degree-zero invariants the inertia cochains have just the singular degree-one space with the cardinality twist. These calculations can be taken first over the fixed local groups at compact precisions before inverting 2, with inflation using lattice invariants, and the descriptions hold also on characteristic-zero field tests over \(A\). Thus \(U_v\) includes local \(H^0\) and injects in \(H^1\).

For the exact complementary conditions on conjugate duals, cup on unramified restrictions factors through unramified cochains (null for the invariant map), and the orthogonality quasi-isomorphism can be checked on residue fields using the displayed dimensions, local duality and the degree-zero invariants. For example the unramified \(H^1\) dimension is the full invariant \(H^0\) dimension; local Euler characteristic is zero. The product of the two singular determinants over \(q\in S^\circ\) is up to units
\[
                   D_q=P_q(\chi(\mathrm{Fr}_v))P_q(\chi(\mathrm{Fr}_v)^{-1}),
\]
nonzero here by the Frobenius-exponent choice. For multiplicative reduction the coinvariant line has arithmetic Frobenius \(a_q=\pm1\); for additive reduction there is no coinvariant line in characteristic zero.
\end{proof}

\begin{lemma}[Projected Kummer lines]\label{split-local:dyadic-lines}
At each \(v=w,\bar w\) there is a projected Kummer condition \(U_v\), a split free line in \(H^1(K_v,M)[-1]\) (placed in degree one) mapping into \(R\Gamma(K_v,M)\), over \(A\). Here \(H^1(K_v,M)\) and the full complex refer to the limiting local model; the latter is itself split free of rank two in degree one only. The lines are exactly self-annihilating with conjugate transport and Weil pairing. The same assertion holds for a fixed characteristic-zero Hecke plane of a modular abelian variety, after allowing fixed denominators for the coefficient projector and polarization. It also holds with several conductor variables, provided at least one of their dyadic Frobenius exponents is nonzero. The high finite character must remain nontrivial on inertia after every fixed extension used to obtain semiabelian reduction.
\end{lemma}

\begin{proof}
We construct the lines and prove exact orthogonality before defining their logarithmic coordinates.

\smallskip\noindent\emph{The local Shapiro tower.}

Put \(F_v=K_v=\mathbb Q_2\), and let \(D_j/F_v\) here denote the unramified extension of degree \(2^j\) (field notation at \(v\)). In the present case the local \(\theta\)-field \(P_v/F_v\) is cyclic totally ramified. The \(u\)-action is unramified with Frobenius value \((1+u)^b\) for \(v_2(b)=1\). Use first the Shapiro scalar
\[
          \Theta[\operatorname{Gal}(P_v/F_v)]
                        [[\operatorname{Gal}(D_\infty/F_v)]]
\]
with tautological character, \(1+s_0\) the Frobenius generator in the unramified factor. Local cohomology is inverse cohomology over \(H_j^{\rm loc}=P_vD_j\) with corestriction. Project the finite factor by \([\sigma]\mapsto\theta(\sigma)\) (idempotent splitting after inverting 2) and substitute \(1+s_0=(1+u)^b\), a finite flat change by Weierstrass or a cyclic-generator change. This computes exactly the desired diagram after these changes: fixed local groups/actions are available by stabilization at each finite precision, and Shapiro, free coefficient models and the fixed projector denominator apply there and in compact inverse limits.

More generally for several conductor variables unramified at \(v\), split \(\theta\) if needed locally as a totally ramified cyclic-field character \(\theta_P\) (on \(P_v\)) times a finite unramified \(\theta_U\), by the abelian local description. Use the finite factor for \(\theta_P\), and substitute for \(1+s_0\) the product of \(\theta_U(\mathrm{Fr}_v)\) and the conductor-variable Frobenius powers. Provided some exponent is nonzero, the change is again flat, finite flat after adjoining the other variables.

\smallskip\noindent\emph{Uniform bounds for logarithmic lattices.}
At local finite levels the integral Kummer sequence on \(H^1\) with Tate coefficients has ends \(\mathcal A(H_j^{\rm loc})^\wedge_2\) and \(\operatorname{Hom}_{\mathbb Z_2}(\mathcal A^\vee(H_j^{\rm loc})^\wedge_2,\mathbb Z_2)\) for the abelian variety \(\mathcal A=E\), or on full modules before a newform projection; identify via polarization allowing fixed denominators for other varieties. Transitions are norm and dual restriction. Compactness permits exact limits. After finite-character projection the log and dual lattices compare to corresponding additive lattices in the Lie spaces and their trace duals up to \textbf{bounded} 2-powers uniformly in \(j\), and projected log torsion kernel has bounded exponent (one can use the unnormalized projector). Both \(\theta_P^{\pm1}\) have this property. Here and below, ``bounded'' means independent of the unramified degree \(2^j\). The constants may depend on the fixed character \(\theta\), its ramified field \(P_v\), the fixed extension used for semiabelian reduction, the abelian variety, and its coefficient projector. Uniformity as \(\theta\) varies is not needed for these characteristic-zero tests.

Indeed a uniform small lattice exponentiates by fixed ramification, using a model first over \(P_v\), then its unramified extensions.

For the upper and torsion bounds extend to the compositum with a fixed Galois \(L/F_v\) of semiabelian reduction. Component exponents there are bounded since after \(LP_v\) the extensions are unramified. An inertia element \(\nu\) over \(L\), chosen with \(\theta_P(\nu)\ne1\), acts trivially on geometric points of the special identity component there.

To see this last fact one can use the special-fiber maps from the Néron model over \(L\). On identity components they are surjective on geometric points: on auxiliary odd-adic Tate modules rationally they identify inertia invariants (Néron reduction over maximal unramified fields is surjective with odd-divisible formal kernel without odd torsion, and components are finite), unchanged under finite base extension by semistable unipotent inertia. Both components are semiabelian by the semistable reduction criterion, so these Tate images suffice for surjectivity. Thus \(\nu\) indeed fixes the special points by functoriality from \(L\).

After a common component multiple, \((\nu-1)\) of any point therefore enters the formal group, and a uniform further multiple enters the small log range. This proves both upper and torsion bounds on applying the finite projectors (their eigenvalue difference from 1 costs only a fixed denominator). The dual log comparison uses the dual of the torsion-free log lattice, and the different denominator is bounded.

\smallskip\noindent\emph{Normal bases and free local cohomology.}
Now \(\mathcal O_{H_j^{\rm loc}}=\mathcal O_{P_v}\otimes_{\mathbb Z_2}\mathcal O_{D_j}\). Use trace-compatible integral normal bases on the unramified tower. More explicitly, choose \(\beta_0=1\) and choose successively \(\beta_j\in\mathcal O_{D_j}\) with
\[
 \operatorname{Tr}_{D_j/D_{j-1}}(\beta_j)=\beta_{j-1}.
\]
The unramified trace is surjective, so this is possible. The total trace of \(\beta_j\) is a unit. Modulo \(2\), the trace operator for the cyclic group of order \(2^j\) is the top power of its generator minus one. Nonzero total trace therefore makes \(\overline\beta_j\) a generator of the residue-field module over the cyclic group algebra. Lifting gives an integral normal basis, and the chosen bases are trace-compatible.

Use also a rational normal basis on the fixed finite factor. The projected norm-limit Kummer terms on both ends therefore become free via log and its dual after inverting \(2\). On the dual end, use the trace-dual lattice; a self-dual normal basis is unnecessary. The transition dual to inclusion is again trace in additive coordinates, and the different of the fixed ramified extension contributes only a fixed denominator. Sandwiching the lattices between fixed multiples is compatible with these limits by compactness.

Degree zero of the full cohomology upstairs vanishes, and the projected degree-two inverse cohomology vanishes after inverting 2 by local duality and the uniform projected torsion bounds. For an elliptic plane the two free ends have rank one each. The same holds on each Hecke-field plane and differential embedding for its modular abelian variety, by characteristic-zero projection (all projector, isogeny and polarization losses there are fixed). Thus \(H^1\) is itself split free with a direct Kummer line: its two free rank-one ends form an exact sequence, and the quotient is projective.

\smallskip\noindent\emph{Exact orthogonality.}
Cup on conjugately paired Kummer lines vanishes, by Shapiro on finite group quotients using ordinary Kummer cups of translates, then on the inverse limits and under the scalar changes. Hence by ranks and perfect local duality this gives exact orthogonality, also on field tests. The restrictions from the split degree-one lines admit isotropy nullhomotopies, so this is an assertion about the local complexes and their comparison triangles, rather than just the dimensions of their generic cohomology.
\end{proof}

\begin{lemma}[Logarithmic coordinate and finite-level transfer]\label{split-local:log-transfer}
The scalar-extended Kummer line at \(w\) has an \(A'\)-linear isomorphism \(\lambda_w\) given by weighted logarithms. At finite cyclic quotients of the conductor variables, this line specializes to the usual projected Kummer condition. A sequence of corestricted twisted classes satisfying the finite-level Kummer conditions to growing precision therefore belongs to the limiting Kummer line. For classes arising from global points, \(\lambda_w\) is the unnormalized character-weighted orbit sum of their logarithms.
\end{lemma}

\begin{proof}
For \(E\) use \(\log=\log_{\omega_E}\); more generally on a Hecke plane take a nonzero corresponding differential defined over the coefficient extension. With evaluations in the fixed embedding, \(\lambda_w\) is the line functional which, before the Frobenius-variable substitution, takes a Shapiro Kummer system at \(w\) to the group-ring sums
\[
 y_j\longmapsto
 \sum_{\sigma\in\operatorname{Gal}(P_wD_j/F_w)}
       \theta_P(\sigma)(1+s_0)^{k(\sigma)}\log(\sigma y_j)
                              \pmod{(1+s_0)^{2^j}-1},
\]
where \(k(\sigma)\) is the unramified exponent. Extend linearly on coefficients. The convention comes from restricting the character corestriction back up; group multiplication on the Shapiro scalar sends the point-system by the inverse Galois action. These sums respect trace and kill the wrong finite-character components. They are bounded after fixed denominators by the lattice comparison; unramified values may use our completed unramified constants.

This gives the line isomorphism claimed: for example a small-log generator before projection uses a fixed small factor times the product of a rational normal generator for \(P_w\) and the above unramified normal generators. Its evaluation is a nonzero constant times a power-series unit. Indeed the finite cyclic normal resolvent at the embedding is nonzero: after scalar extension, the circulant matrix of a rational normal basis is invertible, and its determinant is the product of the character resolvents. Conjugate coordinates are permuted up to units. The unramified group sum has augmentation \(\operatorname{Tr}_{D_j/F_w}(\beta_j)=1\) at every finite level, so its inverse limit is a power-series unit.

\smallskip\noindent\emph{Specialization and membership.}
We explain how this comparison detects membership and logarithms on transferred arithmetic classes. At simultaneous cyclic group quotients of order \(2^k\) in the conductor variables with at least one Frobenius exponent nonzero at 2, the unramified substitution uses cyclic order \(2^{k-\min v_2(b_j)}\) for large \(k\), \(b_j\) the exponents. It is from the corresponding group ring by free scalar extension there (coset bases, and rescaling by the fixed \(\theta_U\) factor if present). The Kummer line specializes, after inverting 2, exactly to the corresponding projected Kummer conditions at those finite levels: use the trace-compatible small-log generators above and ordinary Kummer dimensions under Shapiro.

Suppose as here the stage base ring class fields for \(2^m\) and the tame conductor together locally contain \(P_vD_n\) for every fixed \(n\) on large stage sets. This follows from containing the \(\theta\)-field and the increasing unramified layers by the Frobenius condition. A corestricted twisted class locally Kummer at the stages up to growing precision then lies in \(U_v\) in the test. Indeed at each fixed cyclic quotient its localization is obtained by first norming to these fixed fields by transitivity (and summing the several local terms, with group translates/weights). Thus it specializes into finite-level projected Kummer by the fixed-group comparisons, closedness and congruences at growing coefficient precision there.

Vanishing in the complementary quotient of \(H^1\) is detected on all such finite group quotients (bounded power series, clearing fixed denominators).

\smallskip\noindent\emph{Compatibility with point logarithms.}
For an un-differenced class from global points this also computes its \(\lambda_w\)-value by the unnormalized full orbit sum of logs with character weights. Grouping as above uses actual trace logs in a fixed local extension at each fixed quotient; Kummer congruences at growing precision imply log congruences there, so even the limit into \(\mathcal V^*\) with fixed extensions agrees with the indicated evaluation. Every projection used here has a fixed denominator. The orbit sums are unnormalized throughout, including as the cyclic quotients increase.
\end{proof}

\begin{lemma}[The unramified center]\label{split-local:unramified-center}
Set \(t=0\) and localize, after inverting \(2\), at the tame center \((u)\). At both dyadic places the full local complex again has a direct Kummer line and a complementary quotient line in degree one. They specialize to the ordinary Kummer condition and its quotient at \(u=0\). After the same measure-side scalar extension used for \(\lambda_w\), the weighted logarithm is a line isomorphism over the center local ring, and the transfer and membership assertions of Lemma~\ref{split-local:log-transfer} hold there. These assertions require no ramified-character hypothesis.
\end{lemma}

\begin{proof}
Use just the unramified Shapiro tower. A uniform polynomial in its Frobenius action with value nonzero at 1, times a fixed scalar, now bounds the log comparison in both directions and kills the possibly unbounded reduction torsion. Indeed first use the component order, which is bounded under unramified Néron base change. On the special identity component, the abelian Frobenius polynomial annihilates the abelian part; its roots have complex absolute value \(\sqrt2\). For a torus split over a residue extension of degree \(e\), one may use \(F^e-2^e\); the unipotent part is killed by a fixed power of \(2\). For an elliptic curve these are respectively the good-reduction Frobenius polynomial, the toric polynomial, and the scalar \(2\) in the additive case. Multiplying the relevant factors and, if necessary, another fixed scalar gives a polynomial \(Q(F)\) with \(Q(1)\ne0\). None of the Frobenius polynomial factors has a root-of-unity zero. This enters the formal group, and a further fixed scalar enters the small formal log lattice at all levels.

Thus the map from a fixed small additive lattice by exponentiation, compatible with norm and inclusion, has cokernel killed by such a bound, as does the dual comparison with the opposite action; degree-two projected cohomology is controlled by the same torsion bound. For general varieties likewise use the abelian Frobenius polynomial, the torus Frobenius relation and a characteristic power on the unipotent part.

After localizing (Frobenius corresponds to the unramified group generator or its inverse) we again have exact Kummer/quotient lines on the plane, specializing to ordinary Kummer and quotient at the center by small-log normal generators. Weighted log is a line isomorphism there, and finite-level transfer and membership work as above after first multiplying by the Frobenius polynomial bound (which is invertible at the center). All the bounds canceled in this argument are units at the characteristic-zero center. The assertion is consequently local near that center; it does not assert an unramified logarithmic isomorphism over the entire power-series ring.
\end{proof}


\subsection{The inert derivative identities}

Here is the derivative used on the characteristic-zero test just described. We detail the point and descent part in a form we can also use on a weight-two primitive Hecke plane with trivial nebentypus. There work first with points and full two-Tate coefficients on an actual modular abelian factor \(\mathcal A\) for the newform orbit, then make coefficient projections in characteristic zero. Parametrization maps, Hecke order choices, polarizations or other fixed changes to commensurable lattices can use fixed denominators, always with a \emph{common} scale for the points being compared. For instance we can use an actual Hecke-stable abelian subvariety and multiply the orbit projector into it by an integer. Write \(a_\ell\) for the good Hecke endomorphism/scalar. One can compute with the integral good Hecke action before making further isogeny/lattice changes. For our current problem simply take \(\mathcal A=E\).

Use here lower base fields \(H_{b_i}\), \(b_i=2^m r_i\), base points as in \eqref{eq:M2} with one fixed orientation and original primitive level, and scalar weights \(\chi\) factoring through these fields at each working precision. The point construction works also for base conductor one and more generally the products of a fixed dyadic conductor with prime-to-level moving tame conductors as above, with fields of units just signs (i.e. \(\mathcal O_K^\times=\{\pm1\}\)), keeping any indicated other fixed Kummer places away from the moving conductors. Add finitely many sequences of fresh rational primes \(\ell=\ell_i\) inert in \(K\), avoiding level, conductors and all already used primes at the stages. Take local Frobenius lifts over \(\mathbb Q\) denoted \(\gamma=\gamma_i\). In each sequence require
\[
 \ell\longrightarrow -1,\qquad a_\ell\longrightarrow0,\qquad
 \rho(\gamma)\longrightarrow J,\qquad J^2=1 ,
\]
at dyadic precisions in the ultrafilter, using the \textbf{full} Tate of \(\mathcal A\) for the divisibilities and Frobenius-square closeness. (We also write \(J\) for the induced matrix on a plane in use; it may depend on the sequence of primes.) We will in particular choose \(J=\rho(\tau)\) for complex conjugation \(\tau\) in the present problem.

We can slow working precisions so that \(2^{a+1}\mid\ell+1\), \(a_\ell\) is divisible by \(2^a\) as an endomorphism, and \(\gamma^2\) is trivial on \(\mathcal A[2^a]\), with \(a\) cofinal. In the complex-conjugation approximation on the full newform factor the trace condition holds on all the Hecke field components (complex conjugation is odd and an involution). High enough divisibility in the maximal order at 2 gives the stated divisibility in the good Hecke order by a fixed index bound.

\begin{lemma}[The finite and transverse planes]\label{split-local:inert-planes}
At an auxiliary inert prime satisfying the preceding full-Tate approximation conditions, the limiting local complex has split free cohomology
\[
 M,\qquad M\oplus M,\qquad M
 \quad\text{in degrees }0,1,2.
\]
Here \(M\) denotes the underlying coefficient plane. The finite and transverse conditions have complementary pure degree-one coordinates, both include degree zero, and both omit degree two. They are exactly self-annihilating for conjugate duality. Their mixed pairing is, up to a unit, \(\langle x,Jy\rangle\), where \(\langle\ ,\ \rangle\) is the alternating Tate pairing. The corresponding lower and upper conditions have the exact complementary dualities required for the square switch.
\end{lemma}

\begin{proof}
Write \(f_{\rm loc}(z)=z(\gamma^2)\), \(s_{\rm loc}(z)=z(\sigma_\ell)\) in limiting degree one at this new place of \(K\), for an inertia element \(\sigma_\ell\) generating tame 2-inertia (in its pro-2 quotient) and the whole relative ring-class inertia of order \(\ell+1\) to be used at \(\ell\). The prime \((\ell)\) splits completely in the prime-to-\(\ell\) lower conductor fields by the principal ideal/order criterion, so the base scalar character is trivial there.

Also \(\gamma^2\) acts trivially on \emph{all} the ring class fields of \(K\), even at upper conductor, by generalized dihedral action.

In the limit the full local complex is split with free cohomology \(M,M\oplus M,M\) in degrees 0,1,2, with the indicated degree-one coordinates (here \(M\) in the local coordinates means the underlying coefficient plane). Indeed discard prime-to-2 inertia kernel as in the local calculations; action is trivial locally modulo increasing precision and the tame conjugating power \(\ell^2\) tends to 1. Thus the inertia/residue cohomology or procyclic resolutions give the free coordinates, with zero differentials in the limit. Unramified inflation is the finite plane (\(s_{\rm loc}=0\)); the transverse plane uses \(f_{\rm loc}=0\), including degree zero and omitting degree two as for the finite plane. They fit the lower/upper comparisons in the square switch.

In particular under conjugate transport the pure planes are self-annihilating, and the cross pairing is, up to a unit, \(\langle x,J y\rangle\) by the Tate alternating plane form. Indeed the pure scalar cups are zero as in the local calculation following the square switch, and the mixed scalar cup is a unit. Transport at the place can use \(\gamma\) instead of global conjugation by the inner identification on the cochains; it fixes the Frobenius argument, and changes the tame argument by the \(\ell^{\pm1}\)-power; on coefficients there is the Tate transport with possible scalar character unit. This yields exact complementary dualities also for lower/upper and their triangles (with split local comparisons); for untwisted coefficients one also has the ordinary nonconjugated version with pairing \(\langle x,y\rangle\). For \(T_2 E\) use the unimodular integral plane form, and after inverting 2 any fixed constant plane-form normalization on a characteristic-zero polarized Hecke component costs only units in horizontal valuations. The local identity itself below does not require inverting \(2\) in the Tate limit.
\end{proof}

\begin{lemma}[CM derivatives and uniform descent]\label{split-local:descent}
For a finite collection of auxiliary inert prime sequences satisfying the full-Tate approximation conditions, choose the CM points compatibly over their conductor-product ring class fields. After multiplication by one common nonzero integer \(L\), their derivative Kummer classes descend to the lower conductor fields to arbitrarily high retained \(2\)-power precision. They satisfy Kummer conditions at the fixed places over \(2N\), and unramified conditions at the fixed odd places after a further common component multiple. The same construction applies to a fixed modular abelian factor before characteristic-zero coefficient projection.
\end{lemma}

\begin{proof}
\smallskip\noindent\emph{The conductor-product system.}
For a subset \(I\) of switched primes at a stage, use CM points with conductor \(b_i\prod_{\ell\in I}\ell\), obtained by simultaneous product quotients of the base point by one descending isogeny at each \(\ell\), transporting the level. These choices can extend the previous ones when adding a switch. They are defined over the corresponding upper ring class fields by CM reciprocity: the order changes as stated at the new inert primes, transported level still has stabilizer containing the completed order units. The relative groups over \(H_{b_i}\) are products of independent cyclic inertia groups of orders \(\ell+1\); this is the order class exact sequence and Chinese remaindering, using the unit hypothesis. In particular choose the above generators independently from inertia, acting trivially on the fields with \(\ell\) omitted.

At an upper and omit-\(\ell\) lower pair of points \(Y,X\) respectively in \(\mathcal A\), one has
\[
                 {\rm Tr}_\ell Y=a_\ell X,\qquad
                 \widetilde Y=F_\ell\widetilde X ,
\]
where the first formula if needed is after killing fixed cuspidal torsion, tilde is good reduction at the chosen prime, and \(F_\ell\) is arithmetic residue Frobenius.

These identities work for all matching translates. The trace traverses all the cyclic lines in the good Hecke correspondence, since \(\ell\) is inert and prime to lower conductor. Reduction below is supersingular; all the degree-\(\ell\) quotients there reduce via Frobenius with the transported prime-to-\(\ell\) data. This gives the second relation at the modular point and hence on the abelian image (avoid any fixed primes of bad models/maps). Use maps based at a rational cusp. Also \(F_\ell^2-a_\ell F_\ell+\ell=0\) on the abelian factor of good reduction by the good Eichler--Shimura relation (it can be checked on the rational Tate realization).

\smallskip\noindent\emph{Derivative operators and global descent.}
Use
\[
             D_I=\prod_{\ell\in I}D_\ell,\qquad
             D_\ell=\sum_{j=1}^{\ell} j\sigma_\ell^j .
\]
The Kummer of the upper point acted on by \(D_I\), modulo \(2^a\) after a common cusp/point multiplier, is invariant relatively by \((\sigma_\ell-1)D_\ell=\ell+1-{\rm Tr}_\ell\) and the divisibilities.

A further uniform multiple descends it to \(H_{b_i}\) in cohomology. Indeed \(\mathcal A[2^\infty]\) invariants over these upper fields have bounded exponent. Take an element \(h\in G_K\) acting by a near-one scalar with square not one on \(T_2\mathcal A\) by Bogomolov's homothety theorem for an abelian variety over a number field~\cite{Bogomolov80}. If the scalar is \(c\), then \(h\tau h\tau\) acts by \(c^2\) on the full Tate module. It fixes every ring class field: the image of \(\tau h\tau\) is the inverse of the image of \(h\) in every ring class quotient. Consequently \(c^2-1\) kills all the \(2\)-power torsion defined over any of the upper fields. Its valuation is fixed, giving the required uniform exponent bound.

In inflation--restriction, both the obstruction to descending a relatively invariant Kummer class and the ambiguity of its descent have coefficients in these torsion invariants. They are therefore killed by a fixed power of \(2\). Multiplying by that power gives descent; reducing the coefficient precision by a further fixed amount removes its ambiguity. The resulting retained precision still tends to infinity. These losses are independent of the moving primes and are absorbed into the common point multiplier.

\smallskip\noindent\emph{Descent at the fixed places.}
At \(2N\) the upper extension for these derivatives is unramified relative to the base, with base component exponents uniformly bounded. In fact the ring class fields for the conductor products are composita over the Hilbert field by the same order class description; at 2 the base is thus unramified above a fixed ramification field (fixed \(m\)), and at fixed prime-to-conductor odd support the base itself is unramified over \(K\). Use unramified Néron base change.

The kernel on local Weil--Châtelet in unramified descent is killed by the component group order: for connected points the cyclic unramified \(H^1\) vanishes by Lang on the reduction and the induced additive successive quotients on the formal kernel (unramified residue normal bases, then passage up the complete filtration). Thus a uniform additional multiple ensures the local Kummer condition modulo \(2^a\) on the base there. At fixed odd places as above another component multiple gives unramified Kummer by 2-divisibility of connected points over the maximal unramified field and lifting.

At good places outside allowed support and all derivative primes, the descended class is already unramified since inertia doesn't change going upstairs and prime-to-residue Kummer there is unramified. Use a common multiplier \(L\ne0\) incorporating the fixed losses. Across further switches the bounds allow common scales. If a scale is enlarged, it is enlarged on all classes in the same comparison. Thus each matched derivative relation uses the same \(L\) on its two sides. For a fixed finite sequence of switches, all losses can be incorporated in one such nonzero \(L\).
\end{proof}

After descent, weight-corestrict to \(K\) with \(\chi\), and transfer by the finite models, giving classes \(Z_I\). The weights are the unnormalized character-corestriction weights fixed above.

\begin{proposition}[Inert derivative identity]\label{split-local:derivative}
Fix a finite collection of auxiliary inert prime sequences satisfying
the full-Tate approximation conditions above
Lemma~\ref{split-local:inert-planes}.  Over the lower conductor fields
\(H_{b_i}\), form the compatible CM points and derivative operators
\(D_I\) of Lemma~\ref{split-local:descent}, using its common nonzero
multiplier \(L\) for this collection.  Let \(Z_I\) be their descended,
\(\chi\)-weighted corestrictions to \(K\), transferred to the limiting
coefficient model.  In the finite and singular coordinates
\(f_{\rm loc},s_{\rm loc}\) of
Lemma~\ref{split-local:inert-planes}, at a newly switched prime
\(\ell\notin I\) one has
\[
 s_{\rm loc}(Z_{I\cup\{\ell\}})=J f_{\rm loc}(Z_I),
                 \qquad f_{\rm loc}(Z_{I\cup\{\ell\}})=0,             \tag{K}\label{eq:K}
\]
and the transverse condition in the second identity holds at every switched prime. The identities persist after any fixed equivariant projection onto a Hecke plane, with its denominators absorbed in the common \(L\).
\end{proposition}

\begin{proof}
\smallskip\noindent\emph{The singular coordinate.}
Check the identities before corestriction. Apply the other derivatives to \(Y,X\) first and keep the notation; the matching trace and reduction relations still hold. Choose divisions
\[
                  2^a Q=L X,\qquad 2^a Z=L D_\ell Y
\]
so the respective descended cocycles restrict to their Kummer coboundary expressions upstairs (absorb changes by a torsion coboundary). Thus the finite evaluation below at \(\gamma^2\) reduces to \((F_\ell^2-1)\widetilde Q\). Write \(b\) for the upper descended cocycle on \(G_{H_{b_i}}\). Then
\[
                  U=(\sigma_\ell-1)Z-b(\sigma_\ell)
\]
is fixed by the upper-field subgroup. Indeed, for an element \(g\) of that subgroup, normality gives \(\sigma_\ell^{-1}g\sigma_\ell\) in the same subgroup, and the cocycle identities together with \(b(g)=(g-1)Z\) yield \(gU=U\).

Multiplication by \(2^a\) gives
\[
 2^a U=L\big((\ell+1)Y-a_\ell X\big).
\]
The endomorphism divisibilities ensure that
\[
 L\left(\frac{\ell+1}{2^a}Y-\frac{a_\ell}{2^a}X\right)
\]
is itself an upper-field point. Its difference from \(U\) is \(2^a\)-torsion defined over that field, whose exponent is bounded by Lemma~\ref{split-local:descent}. Inertia acts trivially in reduction. Hence \(b(\sigma_\ell)\) reduces, ignoring bounded-exponent torsion, to
\[
      (a_\ell-(\ell+1)F_\ell)\widetilde Q
                         =(F_\ell-a_\ell)(F_\ell^2-1)\widetilde Q.
\]
The equality follows directly from the Eichler--Shimura polynomial:
\[
 (F_\ell-a_\ell)(F_\ell^2-1)
 =F_\ell(F_\ell^2-a_\ell F_\ell)-F_\ell+a_\ell
 =a_\ell-(\ell+1)F_\ell.
\]
Projection to the retained precision kills the bounded-exponent errors. Prime-to-\(\ell\) torsion injects under good reduction, so the reduction computation determines the Tate coordinate. Since \(F_\ell-a_\ell\) tends to \(J\), it gives the first identity in \eqref{eq:K}.

\smallskip\noindent\emph{The transverse condition.}
For the other identity \(b(\gamma^2)\) reduces to \((F_\ell^2-1)\widetilde Z\), and
\(\widetilde{L D_\ell Y}=L\ell(\ell+1)F_\ell\widetilde X/2\), with \(F_\ell^2\widetilde X=\widetilde X\). Because \(2^{a+1}\mid\ell+1\), an \(F_\ell^2\)-fixed \(2^a\)-division of this reduced point is explicitly
\[
 L\frac{\ell(\ell+1)}{2^{a+1}}F_\ell\widetilde X.
\]
Any other division differs from it by \(2^a\)-torsion. The imposed Frobenius-square condition makes \(F_\ell^2-1\) kill that ambiguity to the retained precision. Hence the finite coordinate of the derived class vanishes.

These calculations work on every translate from the base group over \(K\) (commuting on the points with the relative action), before summing with character weights. Local actions make the coboundary choices in the coordinates immaterial to retained precision. This proves \eqref{eq:K}. Repeating the calculation at an earlier switched prime, with all the other derivative operators already applied, gives the same vanishing there. A fixed equivariant change to a projected coefficient commutes with the calculation; its denominators are cleared in the common \(L\).

Old cochains/classes in successive comparisons can always be used by unramified inflation and reduction to slower precisions; \eqref{eq:K} only uses the matched restrictions of the descended classes.
\end{proof}

\begin{corollary}[Selmer membership and the logarithmic comparison]\label{split-local:point-measure}
For the horizontal test currently considered the \(Z_I\) over \(A\) lie in the Selmer problems with \(U_v\) at \(2N\), full at \(r\), and the switched planes as indicated (finite at additional unswitched derivative places). For \(Z=Z_\varnothing\), their logarithmic comparison with the paired measure is
\[
 B_r(\theta,u)=\epsilon\,
 \lambda_w(\operatorname{loc}_wZ)^2\prod_{q\in S^\circ}D_q,
 \qquad \epsilon\in(A')^\times.\tag{B1}\label{eq:B1}
\]
\end{corollary}

\begin{proof}
Corestriction preserves the unramification just checked away from dyadic and moving support. At odd bad places the inertia restriction class vanishes at every compact precision, hence in the local model (inertia cochains here can be taken on the fixed group; no inverse-limit obstruction from the finite invariants); use inflation.

At 2 use the finite-level Kummer membership check in the preceding local comparisons. Thus the global classes lift, uniquely in degree one since the maps of the conditions include \(H^0\) and inject in degree one.

For \(Z=Z_\varnothing\), the logarithm comparison of Lemma~\ref{split-local:log-transfer}, together with \eqref{eq:M2} in both orientations, gives \eqref{eq:B1}.

The Shapiro weights evaluate the full \(L\)-scaled class sums, the opposite primitive-level orientation differs by Atkin--Lehner and a translation, and the Gauss sums, \(c_E,L\) are fixed nonzero constants with 2 inverted. The two Euler translations at each odd prime multiply exactly with the conjugate arguments. The bounded logarithmic lattice comparisons justify all these limits while retaining the unnormalized sums over the varying class groups.
\end{proof}


\subsection{Horizontal divisibility and the integral unit quotient}

\begin{proposition}[The horizontal determinant inequality]\label{split-local:horizontal}
For a sufficiently high finite character \(\theta\), suppose \(B_r(\theta,u)\ne0\). Impose the unramified conditions \(U_v\) at the fixed odd places and the Kummer conditions at both dyadic places, and call the resulting Selmer complex \(C_F\). Let \(Z=Z_\varnothing\) and write \(\operatorname{loc}_wZ=p_0e\) for a basis \(e\) of the Kummer line over \(A\). Put
\[
 L_{\rm prim}=\frac{L_{\rm alg,r}(\theta,u)}{\prod_{q\in S^\circ}D_q}.
\]
At every height-one valuation \(v\) of \(A\), the complex has generic cohomology of rank one in each of degrees one and two, and
\[
 d(C_F,(Z,Z^\vee))=2v(p_0)-v(L_{\rm prim})\ge0.
 \tag{D2}\label{eq:D2}
\]
The same conclusion holds for the Hecke-plane variant provided its strict determinant is nonzero, its point-measure and local comparisons satisfy the same hypotheses, its restriction to \(G_K\) is absolutely irreducible, and the stated high-character condition holds.
\end{proposition}

\begin{proof}
\smallskip\noindent\emph{The strict comparison.}
Take the high-\(\theta\) test over \(A\) as above; \(L_{\rm alg,r}(\theta,u)\ne0\) by \eqref{eq:A3}. First replace the odd bad full conditions by \(U_v\) there, still using strict at \(w\), full at \(\bar w\); this strict problem is generically acyclic since the removed singular blocks are. Put
\[
       L_{\rm prim}= L_{\rm alg,r}(\theta,u)\Big/\prod_{q\in S^\circ}D_q .
\]
Thus at each height-one DVR of \(A\) the strict problem has \(d(-,1)=-v(L_{\rm prim})\). Now write \(C_F\) for the problem using \(U_v\) also at both dyadic places. If \(B_r(\theta,u)\ne0\), \eqref{eq:B1} gives nonzero localization of \(Z\); writing \(\operatorname{loc}_w Z=p_0 e\) for an \(A\)-basis \(e\) on that Kummer line, \eqref{eq:D} applies by the local comparisons. Thus \(C_F\) has generic ranks one and one, and \eqref{eq:D} gives the equality in \eqref{eq:D2}.

\smallskip\noindent\emph{Detecting two independent fiber classes.}
We prove nonnegativity at every such DVR by the square-switch test of Section~\ref{sec:foundations}, using \eqref{eq:K}. At the generic and DVR-fiber fields all Selmer degree-one spaces used in the square switch inject into unrestricted degree one and thus into abstract crossed classes on \(\mathcal G_K=\prod_i G_K\) by the evaluation lemma. Global invariants are zero and so are the Selmer degree zeros. Exact duality using conjugate transport gives amplitude in degrees 1 and 2; this works also on the paired lower and upper conditions at the extra inert places described above (old diagrams identified via unramified inflation).

In detail let \(k_0\) be the fiber field at this DVR. At a step with old finite condition, degree-one generic rank one and nonzero generic derived point class \(Y=Z_I\), suppose \(H^1\) on the fiber has dimension \(>1\). Take two independent classes there including a primitive reduction from the generic line in the two-term model. Use Shapiro also abstractly into the induction on \(\mathcal G=\mathcal G_K\rtimes\langle\tau\rangle\), with diagonal conjugation. The two \(M\)-planes of the induction over \(\mathcal G_K\) are absolutely simple: already constant \(G_K\) sequences act on Tate via the absolutely irreducible representation (non-CM open image), up to the scalar twists. Their determinants are unequal: the two summands have scalar twists \(\chi\) and \(\chi^{-1}\), and equality of their determinants would force \(\chi^4=1\), contrary to the high dyadic inertia condition. The involution \(\tau\) exchanges these nonisomorphic absolutely simple summands, so the induction is itself absolutely simple. This check works just the same given Tate absolute simplicity over \(K\) on a Hecke plane and the stated high-character condition. For a non-CM elliptic curve the required characteristic-zero absolute simplicity follows from the open image theorem~\cite{Serre72}.

Take the product kernel \(\mathcal H\) of all termwise full Tate and ring class actions within \(\mathcal G_K\), normal also under \(\tau\). Restriction detects the crossed classes since the constant element \(h\tau h\tau\) used in the descent is central modulo \(\mathcal H\) and acts by a scalar whose difference from 1 is invertible here. Thus joint evaluations of the two classes on \(\mathcal H\) span two copies of the induction by the simple-module evaluation test. On \((\tau\eta)^2\) for \(\eta\in\mathcal H\) these evaluations are additive in \(\eta\), given by \(1+\tau\) on the evaluations at \(\eta\). Projection back to the original summand over \(\mathcal G_K\) consequently still spans its two copies. Some \(\eta\) gives rank two there, by the quadratic determinant test.

\smallskip\noindent\emph{A switch preserving the determinant valuation.}
By Chebotarev choose new good inert primes with \(\gamma_i\) agreeing with \(\tau\eta_i\) to the required increasing precisions on the previous evaluation data and full Tate actions. Then \(J=\rho(\tau)\), and in the old problem inflated to the new finite plane the \(f_{\rm loc}\) evaluation is exactly the above test on the fiber. Hence its rank there is two; it is nonzero also generically (use the primitive class) and so tests \(Y\) generically nontrivially.

The new derivative class belongs to the transverse problem and gives exactly \eqref{eq:K}, with previous plane conditions retained. The compatible square switch comparison now preserves the class-functional valuation at this DVR and cuts fiber dimension by two, retaining generic rank one and a nonzero class there. Continue on the same ultrafilter. At each step the fiber dimension drops by two, whereas the nonzero generic class persists, so the dimension remains at least one. After finitely many steps there is no surplus and both degree-one and degree-two fiber dimensions are one.

For completeness, at this terminal DVR the minimal model has one free term in each of degrees one and two. The generic cohomology is also one-dimensional in each degree, so its differential is zero. Exact self-duality identifies these two free cohomology lines as integral duals. The terminal derived class is integral, hence its class-functional tensor has nonnegative valuation. The square-switch equalities carry that valuation back to \((Z,Z^\vee)\). All common fixed-scale multipliers are units in the horizontal DVR, since \(2\) has been inverted. This proves \eqref{eq:D2}.
\end{proof}

\begin{corollary}[Integral unit comparison]\label{split-local:unit}
For the reducible split-product construction,
\[
 B_r=L_{\rm alg,r}\,U_r,
 \qquad U_r\in\mathcal V^*[[t,u]]^\times.
\]
Likewise \(B/L_{\rm alg}(C(t))\) is a unit in \(\mathcal V^*[[t]]\).
\end{corollary}

\begin{proof}

For each sufficiently high \(\theta\) with \(B_r(\theta,u)\ne0\), Proposition~\ref{split-local:horizontal} gives \(p_0^2/L_{\rm prim}\in A\) by the height-one valuation criterion in the regular UFD \(A\). By \eqref{eq:B1} and the logarithmic line isomorphism,
\[
                          B_r(\theta,u)\in L_{\rm alg,r}(\theta,u) A'.
\]
This is automatic if \(B_r(\theta,u)=0\). Weierstrass division in \(u\) using \eqref{eq:A3} and the high-character remainder test therefore prove
\[
                          B_r\in L_{\rm alg,r}\,\mathcal V^*[[t,u]]
\]
integrally. By \eqref{eq:M4} and \eqref{eq:A2} at \(u=0\), cancelling the same nonzero \((a_*-2)^2\), \(B\) is likewise divisible by \(L_{\rm alg}(C(t))\). By \eqref{eq:M5} and \eqref{eq:A1}, the reductions modulo \(2\) of its numerator and denominator have the same finite order in \(t\). Their integral quotient therefore has nonzero constant coefficient modulo \(2\), and is a unit. More explicitly, specializing the two-variable quotient at \(u=0\), relations \eqref{eq:M4} and \eqref{eq:A2} identify it with \(B/L_{\rm alg}(C(t))\) up to an integral unit, after cancellation of the same nonzero factor \((a_*-2)^2\). Thus its specialization at \(u=0\) is a unit. Its constant coefficient at \((t,u)=(0,0)\) is consequently a unit, proving that \(B_r/L_{\rm alg,r}\) is already a unit in the two-variable ring.
\end{proof}

\subsection{The center: rank detection and the exact arithmetic factor}

\begin{proposition}[Rank detection from Selmer corank one]\label{split-local:center-rank}
Under the hypotheses of Proposition~\ref{split-measures:reducible-input}, assume \(s_2(E/K)=1\). The conductor-one Hilbert trace \(P\) in \eqref{eq:M3} is nontorsion. Consequently \(L(E/K,s)\) has a simple zero at \(s=1\).
\end{proposition}

\begin{proof}

We work at \(t=0\), retaining the tame variable \(u\). The nonvanishing of the specialization of a Selmer generator is part of the argument, so no central localization hypothesis is imposed. Localize on that tame line at the \((u)\)-DVR after inverting 2 (and similarly for extended coefficients on the measure side, still testing order in \(u\)). Both series there are nonzero by \eqref{eq:A3} and the unit equality. The odd fixed-place complexes are contractible near this test center: at \(u=0\) they vanish by local duality, the odd-place Euler formula and absence of rational Tate invariants. The added \(r\)-places likewise vanish by construction (\(2-a_*\ne0\)). At \(2\) use Lemma~\ref{split-local:unramified-center}, which supplies the Kummer lines and logarithmic isomorphism at the center DVR. Let \(C_F\) now use these two Kummer conditions, and \(Z(u)\) be the global Kummer class by weighted traces of the conductor-\(r_i\) CM points (disk centers, original level); it belongs to those conditions. The \(t=0\) measure formula following \eqref{eq:M2} gives there
\[
             B_r(0,u)=\text{unit}\cdot\lambda_w(\operatorname{loc}_w Z(u))^2 .
\]
Indeed this is the disk-center logarithm with both orientations as before, retaining additional connected-2 Euler translations, all Euler factors units at this DVR by their values at \(u=0\). For taking the limit of the weighted logs one may first multiply by the polynomial bound of the unramified comparison, evaluated on the group power translating by Frobenius (or its inverse as appropriate). This moves the actual log points into uniformly bounded lattices and specializes to a nonzero scalar at the center. At each fixed cyclic quotient the resulting sums agree in the limit with the line-functional values by corestricting first to the fixed unramified levels, as in the local membership and log checks. Thus one cancels only units of this DVR in asserting the identity.

\smallskip\noindent\emph{The two free central cohomology lines.}
Strict generic acyclicity, the nonzero logarithm, and \eqref{eq:D} again apply. They now give generic ranks one and one for \(C_F\) and \(d(C_F,(Z(u),Z(u)^\vee))=0\), by the unit quotient and the log isomorphism.

The special fiber is also of ranks one and one only: the two lines specialize to usual Kummer; the new allowed \(r\)-places at \(u=t=0\) cause no characteristic-zero change by unramified inflation and local acyclicity. We therefore have the actual ordinary rational Selmer in degree one by fixed-diagram comparison, so use \(s_2(E/K)=1\) and duality. Let \(R\) denote this characteristic-zero \((u)\)-DVR. Its special fiber has just one cohomology line in each of degrees one and two. A minimal free model for \(C_F\) over \(R\) therefore has the form
\[
 [\,R\xrightarrow{\delta}R\,]\qquad\text{in degrees }1,2.
\]
The generic cohomology already has rank one in both degrees, which forces \(\delta=0\). Hence both cohomology groups are free lines over \(R\), paired perfectly by the integral duality over \(R\).

Choose a generator \(y\) of the degree-one line and write \(Z(u)=a(u)y\), with \(a(u)\in R\). The class-functional determinant tensor then has valuation \(2v_u(a(u))\): the two free lines are integral duals and \(y\) has unit determinant volume. Its valuation is zero by the unit comparison. Thus \(a(u)\) is a unit, and
\[
 Z(0)\ne0.
\]
This establishes nonvanishing in global cohomology at the specialization from the Selmer-corank hypothesis.

The specialized class \(Z(0)\) is \((a_*-2)\) times the Kummer of the Hilbert trace \(P\) of \eqref{eq:M3}, up to any common point multiplier used: at the stages trace over the order-\(r_i\) layer gives the good Hecke trace less the two horizontal translates before the Hilbert sum; cuspidal torsion is harmless. These are equalities via inflation of the true old classes at constant coefficients, so pass by the finite diagrams. The factor \(a_*-2\) and the common point multiplier are nonzero in characteristic zero. It follows that \(P\) is nontorsion. Formula~\eqref{eq:GZ-E}, with the split twist sign, now shows that the product has a simple zero~\cite{GZ86,CST14}.
\end{proof}

\begin{proposition}[The exact factor at the center]\label{split-local:exact-factor}
In the reducible split-product setting, assume that \(L(E/K,s)\) has a simple zero at \(1\), either as a hypothesis or by Proposition~\ref{split-local:center-rank}. Then
\[
 X(E)+X(E^{-d})=0.
\]
In the notation of \eqref{eq:Kum}, the exact arithmetic identity producing this equality is
\[
 2(n_P+\tau_g)=2v_2(c_E)+2\sum_\ell v_2(c_\ell)+s_K.
 \tag{E}\label{eq:E}
\]
\end{proposition}

\begin{proof}

Compute at the one-variable center \(C(t)|_{t=0}\). Here we have rank one and finite Tate--Shafarevich over \(K\) by classical Gross--Zagier--Kolyvagin and the quadratic isogeny~\cite{GZ86,Kolyvagin90}, with \(P\) spanning and having nonzero local log.

Impose integral Kummer \(U_v\) at odd places of the support first, keeping strict/full at \(w,\bar w\); denote the resulting complex by \(C_*\). This is rationally acyclic by the converse line comparison of \eqref{eq:D} with all-Kummer. Thus \(C(0)\) is rationally acyclic too by local duality at the added full places.

Use \(v=v_2\), \(n_P,\tau_g,s_K,l,\tau_\ell\) as for \eqref{eq:Kum} and the Haar comparisons (in particular \(s_K=v_2(\#\Sha(E/K))\)). At \(w\) a degree-one free basis with log \(2^l\) has \(d(U_w,e)=\tau_2\). At each odd place over \(q\mid N\), the full-place quotient is the shifted dual of integral Kummer with only degree-two torsion of length \(\tau_q\). Thus
\[
 \begin{split}
 d(C_*,1)
   &=2n_P+2\tau_g-s_K-2\big(v(\log_{\omega_E}P)-l+\tau_2\big),\\
 -v(L_{\rm alg}(C)(0))&=d(C(0),1)=d(C_*,1)-2\sum_{q\in S^\circ}\tau_q,
 \end{split}
\]
by \eqref{eq:D}, \eqref{eq:Kum} and the comparison triangles. The integral unit quotient in Corollary~\ref{split-local:unit} specializes at \(t=0\), so
\[
 v(L_{\rm alg}(C)(0))=v(B(0)).
\]
Formula~\eqref{eq:M3} evaluates the right-hand side as
\[
 v(B(0))=-2v(c_E)+2v(\log_{\omega_E}P)
                  +2\sum_{\ell\mid2N}v(P_\ell(1)).
\]
Combining the displayed determinant comparisons with this equality cancels the two logarithm valuations and yields
\[
 2(n_P+\tau_g)
 =s_K+2v(c_E)+2(\tau_2-l)
       +2\sum_{q\in S^\circ}\tau_q
       -2\sum_{\ell\mid2N}v(P_\ell(1)).
\]
The Haar identities
\[
 \tau_2-l=v(c_2)+v(P_2(1)),\qquad
 \tau_q=v(c_q)+v(P_q(1))\quad(q\ne2)
\]
cancel every Euler factor and give exactly \eqref{eq:E}. The two places of \(K\) above each odd support prime account for the factor \(2\) in the singular determinant contribution; the dyadic Kummer volume retains \(\tau_2-l\).

Finally, \(n_P\) is the index valuation in the full free Mordell--Weil lattice and \(\tau_g\) records its torsion. Therefore
\[
 n_P+\tau_g=v([E(K):\mathbb ZP]),
\]
and \eqref{eq:G} identifies \eqref{eq:E} with \(X(E)+X(E^{-d})=0\). Together with Proposition~\ref{split-local:center-rank}, this proves Proposition~\ref{split-measures:reducible-input}.
\end{proof}


\subsection{A characteristic-zero Heegner rank test}

The next result supplies the low analytic rank Selmer input on the Hecke components used later. Its proof uses Gross--Zagier, the inert derivative identity of Proposition~\ref{split-local:derivative}, and global reciprocity. It is independent of the horizontal determinant inequality, the residual unit comparison, and the rank deduction at an initially unknown center.

\begin{proposition}[Characteristic-zero Heegner rank test]\label{split-local:heegner-rank}
Let \(f\) be a normalized primitive weight-two newform of trivial nebentypus and level \(N_f\), and let \(\mathcal A\) be its modular abelian factor. Let \(K\) have odd fundamental discriminant \(-D<-4\), prime to \(N_f\), and suppose every prime dividing \(N_f\) splits in \(K\). Thus \(\mathcal O_K^\times=\{\pm1\}\), and the split height formula \eqref{eq:GZ} applies. Let \(F_f\) be a finite extension of \(\mathbb Q_2\) containing the image of a fixed Hecke coefficient embedding, and let \(V_f\) be the corresponding two-dimensional Tate representation over \(F_f\), allowing a further finite scalar extension when needed. Assume \(V_f|_{G_K}\) absolutely irreducible. Then
\[
 \begin{gathered}
 L(f/K,s)\ \text{has a simple zero at }1\\
 \Longrightarrow\quad
 H^1_f(K,V_f)=F_f y,\qquad y\ne0,\qquad \tau y=-w(f)y,
                                                                    \tag{K0}\label{eq:K0}
 \end{gathered}
\]
where \(y\) is the Kummer class of the projected Hilbert trace. The Selmer notation on the right means ordinary Kummer finite conditions, with conjugate rational twist components over \(\mathbb Q\) recovered by quadratic isogeny. In particular the summand whose \(L\)-function has nonzero central value has zero such Selmer and the simple-zero summand rank one, since \(w(f)\) is the functional sign over \(\mathbb Q\) and the split product sign is minus. Dual Tate via a polarization gives the same assertions rationally. The conclusion concerns the characteristic-zero Selmer rank at the indicated coefficient. For a general \(\mathcal A\), no assertion of finiteness for its entire prime-primary Tate--Shafarevich group is included.
\end{proposition}

\begin{proof}
\smallskip\noindent\emph{The Heegner class and its sign.}

Formula~\eqref{eq:GZ} gives a nonzero projected Hilbert point up to torsion~\cite{GZ86,CST14}; its Kummer class \(y\) projects nontrivially on \(V_f\). The rational Mordell--Weil space for the newform orbit is a vector space over its Hecke field. A nonzero vector stays nonzero after scalar extension through each embedding of that field; scalar extension to a field is faithful. Kummer injectivity is Hecke-equivariant, so the projected Kummer class is nonzero at the specified \(2\)-adic coefficient embedding.

Conjugation reverses all level orientations in the Hilbert sum. This agrees up to ideal translation with the full Fricke action (divide by the oriented level subgroup and use the dual subgroup); base cusp differences are torsion. The Fricke eigenvalue is \(-w(f)\), giving the asserted sign \(s_y\) for \(y\). 

\smallskip\noindent\emph{Local signs and simultaneous evaluations.}
Apply \eqref{eq:K} at base conductor one, with no scalar twists, working with the usual untwisted Tate pairings and variable inert primes. Use the alternating plane pairing on \(V_f\); up to coefficients this comes from the polarized ordinary pairing and Hecke projection (Rosati fixes the field). The initial problems therefore use exact orthogonal conditions. At fixed places descent gives ordinary Kummer by the proof above (it only needs fixed base local ramification here).

Switched places use the two pure planes; all these conditions are conjugation stable in the limit. At a new place with \(\gamma_i\) tending to complex conjugation on full Tate, the action of global conjugation on localizations of classes is \(J=\rho(\tau)\) in \(f_{\rm loc}\) and \(-J\) in \(s_{\rm loc}\), by the same comparison with \(\gamma_i\) and tame transport. In this proof a ``new prime'' means a fresh auxiliary prime sequence with its limiting arithmetic diagram, as in Proposition~\ref{split-local:derivative}. All finite and singular coordinates, local pairings, and reciprocity comparisons below are taken in those compatible limiting diagrams.

Because complex conjugation is odd on a weight-two Tate plane, each eigenspace of \(J\) is one line. Write \(M_+\) and \(M_-\) for those lines. The alternating pairing vanishes on either line with itself and is nondegenerate between the two lines.

For any nonzero conjugation eigenclasses of sign \(s_x\) in a current Selmer problem, finite evaluation at a new place can be made nonzero simultaneously for finitely many such classes. Use again the product kernel \(\mathcal H\) and cochain evaluations; restriction is injective by the same central element (now even on the plane alone), and the evaluations of each class span the simple plane. On \((\tau\eta)^2,\ \eta\in\mathcal H\), its values are \((1+s_x J)\) of those evaluations, since conjugation comparison is on classes and coboundaries restrict trivially there. Thus they span the corresponding sign line \(M_{s_x}\). Explicitly, on the joint kernel, where the coefficient action is trivial,
\[
 x\big((\tau\eta)^2\big)=(1+s_xJ)x(\eta).
\]
The same formula is valid on the cohomology classes because conjugation comparison holds on classes and coboundaries restrict trivially to that kernel. By additivity in characteristic zero, simultaneous nonzero tests for finitely many classes exist; use Chebotarev as above, retaining all the data including closeness on the full \(\mathcal A\). This works also in the extended diagrams with earlier switches.

\smallskip\noindent\emph{Excluding the opposite sign.}
Suppose an initial nonzero Selmer class \(x\) has sign \(-s_y\). Choose a first prime at which both \(x\) and \(y\) have nonzero finite evaluations. Then
\[
 f_{\rm loc}(y)\in M_{s_y}\setminus\{0\},\qquad
 f_{\rm loc}(x)\in M_{-s_y}\setminus\{0\}.
\]
The derivative of \(y\) (using common nonzero scaling by \(L\) here) has singular evaluation \(J\) times its nonzero finite evaluation. Pair the derivative with \(x\). At the new prime the mixed local cup is nonzero, because its singular vector \(Jf_{\rm loc}(y)\) lies in \(M_{s_y}\setminus\{0\}\), opposite to the nonzero finite vector of \(x\). Every other local cup vanishes by the self-annihilating conditions. Global cup reciprocity says that the sum of these invariants is zero, a contradiction.

\smallskip\noindent\emph{Excluding a second line of the same sign.}
The initial Selmer space is therefore wholly in sign \(s_y\). Suppose its dimension exceeds one, and choose a first prime at which \(y\) evaluates nontrivially. Some nonzero old class \(x\) is in the finite evaluation kernel since the evaluation is in just one \(J\)-line. Its finite and singular coordinates at this prime are both zero, so it lies in the common lower condition and persists in the transverse problem, still in sign \(s_y\). The new derivative \(y'\) has nonzero singular coordinate \(Jf_{\rm loc}(y)\in M_{s_y}\). Conjugation acts as \(-J\) on this coordinate, so the \(-s_y\)-eigenprojection of \(y'\) is nonzero. At a second prime, make the finite evaluations of \(x\) and that eigenprojection nonzero simultaneously. Differentiate the full class \(y'\), using its place in the derivative system. By \eqref{eq:K}, its twofold derivative has singular coordinate \(Jf_{\rm loc}(y')\). The component of this vector in \(M_{-s_y}\) is nonzero and pairs nontrivially with \(f_{\rm loc}(x)\in M_{s_y}\setminus\{0\}\). Any component in \(M_{s_y}\) pairs to zero with \(f_{\rm loc}(x)\), so it cannot cancel the nonzero contribution.

At the first prime, both classes remain in the transverse condition, which is self-annihilating. At all fixed places the Kummer conditions are retained. Thus the second prime gives the only nonzero local cup invariant, again contradicting reciprocity. The Selmer space is exactly the line generated by \(y\), with sign \(s_y=-w(f)\), proving \eqref{eq:K0}.
\end{proof}

\begin{remark}[Use with quadratic companions]\label{split-local:rank-companions}
For a primitive form with an individual nonzero central value, the quadratic-twist nonvanishing results supply a split imaginary companion with a simple zero. For an individual simple zero, they supply such a companion with nonzero central value, under the stated local and sign-compatible prescriptions~\cite{FH95,BFH90}. Proposition~\ref{split-local:heegner-rank} applies whenever the resulting restriction of the Tate plane remains absolutely irreducible.

In particular the latter simplicity hypothesis on the Tate action holds for a CM form when \(K\) is distinct from its inducing quadratic field and the two inducing characters remain distinct on the compositum; it also holds for a characteristic-zero lift of an absolutely irreducible residual plane over \(K\). We will check the relevant bases below. Thus \eqref{eq:K0} supplies these characteristic-zero Selmer rank comparisons directly from the Heegner derivative construction and does not require a residual Euler-system bound.
\end{remark}
